\documentclass[11pt, a4paper]{article}
\usepackage[margin=1in]{geometry}
\usepackage{booktabs}
\usepackage{hyperref}
\usepackage{enumitem}
\usepackage{xcolor}
\usepackage{array}
\usepackage{longtable}
\usepackage{listings}
\usepackage{amssymb}
\usepackage{parskip}
\usepackage{titlesec}

\definecolor{critred}{RGB}{180,0,0}
\definecolor{majblue}{RGB}{0,60,160}
\definecolor{mingreen}{RGB}{0,120,0}
\definecolor{codegray}{RGB}{245,245,245}
\definecolor{darkgray}{RGB}{80,80,80}

\lstset{
  backgroundcolor=\color{codegray},
  basicstyle=\ttfamily\small,
  breaklines=true,
  frame=single,
  rulecolor=\color{darkgray},
  commentstyle=\color{darkgray}\itshape,
  keywordstyle=\color{majblue},
  stringstyle=\color{mingreen},
  xleftmargin=1em,
  xrightmargin=1em,
  aboveskip=0.5em,
  belowskip=0.5em
}

\hypersetup{
  colorlinks=true,
  linkcolor=majblue,
  urlcolor=majblue,
  citecolor=darkgray,
  pdfauthor={Greppmair, Jank, Saffi, Sturgess},
  pdftitle={Revision Plan -- Securities Lending and Information Acquisition}
}

\title{\textbf{Revision Plan}\\[0.5em]
\large Securities Lending and Information Acquisition\\[0.3em]
\normalsize Greppmair, Jank, Saffi, Sturgess}
\date{April 2026}
\author{}

\begin{document}
\maketitle
\tableofcontents
\newpage

% ============================================================
\section{Overview and Journal Target}
% ============================================================

\subsection{Current Status}

The paper presents credible, well-identified evidence that active mutual funds exploit private
information from equity lending operations: lending predicts future position reductions, the
effect is amplified around mandatory short-position disclosures, passive funds show no reaction
(clean placebo), and same-manager lending transmits information more strongly than cross-manager
within-family lending. The identification strategy is sound and the institutional setting is novel
relative to the US literature.

\textcolor{mingreen}{\textbf{Strong JFQA / Review of Finance submission now.}}

\textcolor{majblue}{\textbf{JF borderline: borderline pending four required additions (Section~3).}}

\subsection{Journal Target Recommendation}

\begin{center}
\begin{tabular}{lll}
\toprule
\textbf{Journal} & \textbf{Readiness} & \textbf{Bottleneck} \\
\midrule
JFQA & Ready (post Section~2 fixes) & None blocking \\
Review of Finance & Ready (post Section~2 fixes) & None blocking \\
JFE & 2--4 weeks of additional polish & Sections 3.1, 3.2, 4.1 \\
JF & 2--3 months of additional work & All of Section~3 \\
\bottomrule
\end{tabular}
\end{center}

\subsection{Estimated Revision Timeline}

\begin{itemize}[leftmargin=2em]
  \item \textbf{2--4 weeks}: JFQA/RFS-ready clean draft. Complete all items in Sections~2 and~4.
  \item \textbf{2--3 months}: Full JF package. Complete Sections~2, 3, 4, and 5.
\end{itemize}

% ============================================================
\section{Completed Fixes (This Draft)}
% ============================================================

\textcolor{mingreen}{All items below have been applied to the current working draft.}

\subsection*{Equations}
\begin{itemize}[leftmargin=2em]
  \item[$\checkmark$] \textcolor{mingreen}{Removed phantom $\beta_3$ term from Eq.~(3): \texttt{OnLoan\_{i,t-1} $\times$ Active\_j} appeared in the equation but had no corresponding column in Table~4.}
  \item[$\checkmark$] \textcolor{mingreen}{Removed phantom $\beta_2$ term from Eq.~(4): \texttt{OnLoan $\times$ OnLoan\_{i,t-1}} appeared with no table column or discussion; renumbered $\beta_3 \to \beta_2$, $\beta_4 \to \beta_3$, $\beta_5 \to \beta_4$.}
  \item[$\checkmark$] \textcolor{mingreen}{Applied \texttt{\textbackslash mathit\{\}} to all variable names in math mode: \textit{OnLoan}, \textit{PosChange}, \textit{PosWeight}, \textit{ILLIQ}.}
  \item[$\checkmark$] \textcolor{mingreen}{Replaced \texttt{eqnarray} with \texttt{align} throughout (standard multi-line equation environment).}
  \item[$\checkmark$] \textcolor{mingreen}{Eq.~(1): changed \texttt{Y(\ldots)} notation to \texttt{f(\ldots)} to avoid reusing $Y$ as a function name when $Y$ is already the dependent variable.}
\end{itemize}

\subsection*{Abstract and Introduction}
\begin{itemize}[leftmargin=2em]
  \item[$\checkmark$] \textcolor{mingreen}{"demonstrate" $\to$ "show" in abstract (finance convention).}
  \item[$\checkmark$] \textcolor{mingreen}{"more than doubles" $\to$ "nearly triples" in abstract; matching "nearly triples" in introduction now consistent.}
  \item[$\checkmark$] \textcolor{mingreen}{Introduction timing: "month $t$" $\to$ "month $t-1$" for lending-to-trading lead.}
  \item[$\checkmark$] \textcolor{mingreen}{EFP acronym defined at first occurrence in introduction.}
  \item[$\checkmark$] \textcolor{mingreen}{\texttt{table\_passive\_iv} $\to$ \texttt{table\_passive\_rf} cross-reference corrected.}
\end{itemize}

\subsection*{Tables}
\begin{itemize}[leftmargin=2em]
  \item[$\checkmark$] \textcolor{mingreen}{Table~7 notes: \texttt{\textbackslash scriptsize} $\to$ \texttt{\textbackslash footnotesize}.}
  \item[$\checkmark$] \textcolor{mingreen}{Table~9: added SE Clustering and Sample rows.}
  \item[$\checkmark$] \textcolor{mingreen}{Table~9: added note explaining $N$ discrepancy (39,109 vs 39,003).}
\end{itemize}

\subsection*{Results Section}
\begin{itemize}[leftmargin=2em]
  \item[$\checkmark$] \textcolor{mingreen}{Added explanatory sentence for columns~(5)--(7) of Table~5, which were previously unreferenced in the text.}
\end{itemize}

% ============================================================
\section{Required Before JF Submission (Priority 1)}
% ============================================================

\textcolor{critred}{\textbf{These items are blocking for JF. All must be completed and reviewed before submission.}}

\subsection{Formal IV Design}

\textbf{Issue.} The paper asserts that passive fund lending within the same family serves as an
instrument for active fund lending (endogenous variable: \textit{OnLoan} for active funds).
However, no formal IV table is presented with first-stage F-statistic, instrument diagnostics,
or exclusion restriction discussion. Referees at JF will require this.

\textbf{Expected output.} A new or revised table (Table~6 or Internet Appendix) containing:
\begin{enumerate}[leftmargin=2em]
  \item First-stage regression with coefficient on passive family lending, F-statistic, and
        Kleibergen-Paap rk Wald F-statistic.
  \item Second-stage (2SLS) estimate of position change on instrumented active lending.
  \item A paragraph in the text discussing the exclusion restriction: passive funds mechanically
        hold positions determined by their benchmark, not by information signals, so their
        lending decisions are plausibly exogenous to active fund trading decisions.
\end{enumerate}

\textbf{Stata code template.} See Section~\ref{sec:stata-iv}.

\subsection{Wald Test for Timing Symmetry}

\textbf{Issue.} The paper reports $\hat{\beta}_{\text{lead}} \approx -0.87$ and
$\hat{\beta}_{\text{lag}} \approx -0.30$ (t-statistic $-1.34$), and interprets the asymmetry
as evidence that trading precedes disclosure. However, no formal test of
$H_0: \beta_{\text{lead}} = \beta_{\text{lag}}$ is reported.

\textbf{Expected output.} Report a $\chi^2$ test (or equivalent Wald test from \texttt{lincom})
with p-value. A one-sentence footnote suffices if the test rejects at conventional levels.
If it does not reject, the interpretation of asymmetric timing requires softening.

\textbf{Stata code template.} See Section~\ref{sec:stata-wald}.

\subsection{Return Predictability Table}

\textbf{Issue.} The conclusion states that funds generate gains "at the expense of uninformed
buyers." This welfare claim is unsubstantiated without return evidence. JF referees will demand
either (a) a performance table showing risk-adjusted returns for active lenders in the month
before disclosure, or (b) removal of the welfare language.

\textbf{Expected output.} A table (Internet Appendix or paper body) reporting monthly
Carhart four-factor alphas for:
\begin{itemize}[leftmargin=2em]
  \item Active lending funds vs non-lending active funds in months $[-3, +3]$ relative to
        disclosure events.
  \item Preferably as a panel regression with fund $\times$ time fixed effects and a lender
        indicator interacted with event-month indicators.
\end{itemize}

\textbf{Stata code template.} See Section~\ref{sec:stata-returns}.

\subsection{Mechanical Autocorrelation Diagnostic}

\textbf{Issue.} When a fund reduces its position (sell), it mechanically holds fewer shares
available to lend, so \textit{OnLoan} mechanically declines in the same or subsequent month.
This creates a spurious positive correlation between lagged lending and position changes that
could bias the coefficient of interest. The paper does not address this threat.

\textbf{Expected output.}
\begin{enumerate}[leftmargin=2em]
  \item A robustness row in an existing table excluding lending-termination months (months where
        $\mathit{OnLoan}_{t-1} > 0$ and $\mathit{OnLoan}_t = 0$).
  \item A Wooldridge test for serial correlation in residuals (\texttt{xtserial}), reported in a
        footnote or robustness table.
\end{enumerate}

\textbf{Stata code template.} See Section~\ref{sec:stata-ar}.

\subsection{Remove ``We Plan to \ldots in the Revision'' Language}

\textcolor{critred}{\textbf{Critical: blocks submission at any journal.}}

Three locations contain author notes that were not removed before this review:

\begin{enumerate}[leftmargin=2em]
  \item \texttt{empirical\_strategy.tex}, line 42
  \item \texttt{empirical\_strategy.tex}, line 54
  \item \texttt{information\_channel.tex}, footnote, line 17
\end{enumerate}

Each must be either (a) deleted, or (b) converted to completed-task language (e.g., "We
address this in Section~X" or "Robustness results are reported in Table~X"). No draft with
placeholder author notes should be submitted to any journal.

% ============================================================
\section{High Priority (Priority 2) --- JFQA/RFS Standard}
% ============================================================

\textcolor{majblue}{\textbf{These items should be addressed before submitting to any top field journal.}}

\subsection{Add Adj. $R^2$ to Tables 6, 8, and 9}

Finance panel regressions report Adjusted $R^2$ by field convention. Tables~6, 8, and~9 currently omit it. This requires a Stata re-run adding \texttt{estadd scalar r2\_a = e(r2\_a)} after each \texttt{reghdfe} command.

\subsection{Welfare Language in Conclusion}

The sentence "generating gains for their investors at the expense of uninformed buyers" in the
conclusion makes a strong welfare claim that is not supported by performance evidence in the
paper. Resolution options:

\begin{enumerate}[leftmargin=2em, label=(\alph*)]
  \item Add the return predictability table (Section~3.3) — preferred.
  \item Replace with hedged language: "consistent with funds extracting rents from lending
        counterparties" or similar.
  \item Remove entirely.
\end{enumerate}

\subsection{Weaken ``Consistent Only With'' Language}

\texttt{empirical\_strategy.tex}, line~52 currently reads "consistent only with active
information acquisition." The word "only" is too strong — no formal exclusion test rules
out all alternative mechanisms. Replace with "most consistent with" or "difficult to reconcile
with alternatives other than."

\subsection{Missing Citations}

\begin{enumerate}[leftmargin=2em]
  \item \textbf{Blocher \& Ringgenberg (2023, JFE).} Short interest predicts returns. Directly
        relevant to the return predictability mechanism; a referee will notice its absence.
  \item \textbf{Massa, Qian, Shi, \& Zhang (2015, JF).} Mutual fund managers and information
        acquisition through institutional networks. Closest analogue to the family-spillover
        channel.
\end{enumerate}

\subsection{Fix Descriptive Tables}

\begin{enumerate}[leftmargin=2em]
  \item Replace \texttt{\textbackslash hline} with \texttt{\textbackslash midrule} in Tables~1--3.
        The paper uses \texttt{booktabs} for regression tables; descriptive tables should be
        consistent.
  \item Unify panel label style: Table~3 uses \texttt{\textbackslash textbf\{\}}, Tables~8--9
        use \texttt{\textbackslash textit\{\}}. Standardize to \texttt{\textbackslash textit\{\}}
        throughout.
\end{enumerate}

\subsection{Standardize Table Notes Location}

Currently, tables~4, 6, and~7 have notes embedded in their \texttt{.tex} files, while
tables~5, 8, and~9 have notes in \texttt{main.tex}. Move all notes to their respective
\texttt{.tex} files for maintainability and consistency.

% ============================================================
\section{Suggested Analyses for Stronger Paper}
% ============================================================

\textcolor{darkgray}{Advisory: these additions would strengthen the paper but are not blocking
for JFQA/RFS submission.}

\subsection{Event-Study Figure}

Plot average position changes for active lenders vs.\ non-lenders in months $[-6, +6]$
relative to the mandatory disclosure event. This is the single highest-impact presentation
improvement: a well-executed event study figure conveys the paper's main result visually and
is now standard at JF and JFE.

\subsection{Short-Seller Identity Heterogeneity}

If borrower identity is available in the Bundesbank IFS data or through IHS Markit, test
whether funds respond more strongly to lending when the borrower is a historically accurate
short seller (measured by subsequent price declines). This would substantially strengthen the
information-acquisition mechanism story.

\subsection{SFTR Post-Period Comparison}

The Securities Financing Transactions Regulation (SFTR) was implemented in 2020 and increased
transparency in the lending market. A before/after comparison around SFTR implementation
serves as an additional falsification test: if funds exploit private lending information, the
effect should weaken post-SFTR when that information becomes more public.

\subsection{Multiple Testing Corrections}

Table~9 reports several triple-interaction coefficients. At JF, referees routinely request
Romano-Wolf multiple testing corrections for tables with many simultaneous hypotheses. Apply
\texttt{rwolf} (Romano \& Wolf 2005) to the Table~9 interactions.

% ============================================================
\section{Stata Code Suggestions}
% ============================================================

\textbf{Note:} All code below is a template requiring adaptation to actual variable names,
dataset structure, and fixed-effect specification in use. Variable names (\texttt{onloan\_active},
\texttt{poschange}, etc.) are illustrative placeholders.

\subsection{Formal IV First Stage}\label{sec:stata-iv}

\begin{lstlisting}[language={}]
* ============================================================
* Table 6 - Formal IV: Passive Family Lending as Instrument
* Endogenous variable: active fund lending (onloan_active)
* Instrument: passive fund lending within same fund family
* ============================================================

* --- First Stage ---
reghdfe onloan_active onloan_passive_family ///
    log_aum fund_age expense_ratio turnover ///
    , absorb(securityid#date fundid#date fundid#securityid) ///
    cluster(fundid date)

estadd scalar Fstat = e(F)
estadd scalar r2_a = e(r2_a)

* Kleibergen-Paap rk Wald F-statistic (instrument strength)
test onloan_passive_family
estadd scalar Fstat_excl = r(F)

* Rule of thumb: Fstat_excl > 10 (Stock-Yogo); > 104 for 5% max IV bias

* --- Second Stage (2SLS) ---
* Requires ivreghdfe (ssc install ivreghdfe)
ivreghdfe poschange (onloan_active = onloan_passive_family) ///
    log_aum fund_age expense_ratio turnover ///
    , absorb(securityid#date fundid#date fundid#securityid) ///
    cluster(fundid date)

estadd scalar r2_a = e(r2_a)

* --- Store and output ---
esttab using "Tables/table_iv.tex", ///
    keep(onloan_active onloan_passive_family) ///
    stats(Fstat_excl N r2_a, ///
          labels("F-stat (excl. instrument)" "N" "Adj. R-sq")) ///
    label star(* 0.10 ** 0.05 *** 0.01) ///
    b(4) t(2) replace
\end{lstlisting}

\subsection{Wald Test for Timing Symmetry}\label{sec:stata-wald}

\begin{lstlisting}[language={}]
* ============================================================
* Wald Test: H0: beta_lead = beta_lag
* After estimating the disclosure-timing specification (Table 7)
* ============================================================

* --- Estimate disclosure-timing model ---
reghdfe poschange ///
    c.onloan##i.disclosure_lead   ///   interaction: onloan x lead disclosure
    c.onloan##i.disclosure_contemp ///  interaction: onloan x contemp disclosure
    c.onloan##i.disclosure_lag    ///   interaction: onloan x lag disclosure
    log_aum fund_age expense_ratio turnover ///
    , absorb(securityid#date fundid#date fundid#securityid) ///
    cluster(fundid date)

* --- Wald test: coefficient on lead = coefficient on lag ---
* Method 1: lincom (gives estimate of difference + SE + t-stat)
lincom _b[c.onloan#1.disclosure_lead] - _b[c.onloan#1.disclosure_lag]

* Method 2: test (gives F-statistic and p-value)
test _b[c.onloan#1.disclosure_lead] = _b[c.onloan#1.disclosure_lag]

* Report: chi2(1), p-value in footnote to Table 7
* If p < 0.10: "We reject equality of lead and lag coefficients
*   (chi2(1) = X.XX, p = 0.0XX), consistent with trading preceding disclosure."
\end{lstlisting}

\subsection{Return Predictability Around Disclosure}\label{sec:stata-returns}

\begin{lstlisting}[language={}]
* ============================================================
* Return Predictability: Risk-adjusted returns for lenders
* around mandatory disclosure events
* ============================================================

* --- Step 1: Compute Carhart 4-factor alpha ---
* Merge fund-level monthly returns with factor data (French library)
* carhart_alpha already computed as rolling 36-month window alpha

* --- Step 2: Define event window relative to disclosure ---
* event_month = month relative to disclosure (-6 to +6)
gen event_month = date - disclosure_date  // in months

* Keep [-6, +6] window
keep if event_month >= -6 & event_month <= 6

* --- Step 3: Panel regression with lender x event-month interactions ---
* lender_active = 1 if fund lent the stock in month t-1
reghdfe carhart_alpha ///
    i.lender_active##i.event_month ///
    , absorb(securityid fundid) ///
    cluster(fundid date)

* --- Step 4: Plot coefficients for event-study figure ---
* coefplot (ssc install coefplot):
coefplot, ///
    keep(1.lender_active#*) ///
    vertical ///
    yline(0, lpattern(dash)) xline(6.5, lpattern(dash)) ///
    xtitle("Months relative to disclosure") ///
    ytitle("Abnormal return (%, monthly)") ///
    title("Active lenders vs. non-lenders around disclosure") ///
    saving("Figures/event_study_returns.gph", replace)
graph export "Figures/event_study_returns.pdf", replace
\end{lstlisting}

\subsection{Mechanical Autocorrelation Diagnostic}\label{sec:stata-ar}

\begin{lstlisting}[language={}]
* ============================================================
* Mechanical Autocorrelation Diagnostic
* Threat: lending declines mechanically when positions are sold
* Test: exclude lending-termination months; test AR(1) residuals
* ============================================================

* --- Step 1: Flag lending-termination months ---
* Termination = fund lent in t-1 but not in t
sort fundid securityid date
by fundid securityid: gen term_month = ///
    (onloan[_n-1] > 0 & !missing(onloan[_n-1])) & (onloan == 0)

label var term_month "Lending-termination month (mechanical)"

* --- Step 2: Re-estimate baseline excluding termination months ---
reghdfe poschange onloan ///
    log_aum fund_age expense_ratio ///
    if !term_month ///
    , absorb(securityid#date fundid#date fundid#securityid) ///
    cluster(fundid date)

estadd scalar r2_a = e(r2_a)
* Compare to baseline: if coefficient similar, mechanical story not driving results

* --- Step 3: Wooldridge test for serial correlation ---
* Requires xtserial (ssc install xtserial)
* Run on baseline specification residuals

reghdfe poschange onloan log_aum fund_age expense_ratio ///
    , absorb(securityid#date fundid#date fundid#securityid) ///
    cluster(fundid date) residuals(resid_base)

xtset fundid date
xtserial resid_base
* H0: no first-order serial correlation
* If fail to reject: serial correlation not a concern
* If reject: report AR(1)-robust standard errors as robustness
\end{lstlisting}

\subsection{Multiple Testing Corrections (Table 9)}\label{sec:stata-mt}

\begin{lstlisting}[language={}]
* ============================================================
* Romano-Wolf Multiple Testing Corrections for Table 9
* Applies to triple-interaction coefficients
* Requires: ssc install rwolf
* ============================================================

* --- Estimate all Table 9 specifications and store ---
* (run each reghdfe specification and store in e() or matrix)

* --- Apply Romano-Wolf correction ---
* rwolf syntax: rwolf (models...), reps(B) seed(S) [onesided]
* Note: rwolf requires eststo-stored models

eststo m1: reghdfe poschange c.onloan##i.active##i.hhi_high ///
    [controls], absorb(...) cluster(fundid date)

eststo m2: reghdfe poschange c.onloan##i.active##i.hhi_mid ///
    [controls], absorb(...) cluster(fundid date)

* Romano-Wolf correction across stored models
rwolf (m1: onloan active onloan#active#hhi_high) ///
      (m2: onloan active onloan#active#hhi_mid)  ///
    , reps(500) seed(42)

* Report: original p-values and Romano-Wolf adjusted p-values
* in footnote or table panel
\end{lstlisting}

% ============================================================
\section{Quality Score Summary}
% ============================================================

\begin{center}
\begin{longtable}{lllll}
\toprule
\textbf{Component} & \textbf{Agent} & \textbf{Pre-Fix} & \textbf{Post-Fix} & \textbf{Status} \\
\midrule
\endhead
Spelling/Grammar/Style & Agent 1 & 88/100 & 88/100 & \textcolor{mingreen}{\textbf{PASS}} \\
Internal Consistency & Agent 2 & 78/100 & 82/100 & \textcolor{mingreen}{\textbf{PASS}} \\
Unsupported Claims & Agent 3 & 75/100 & 75/100 & \textcolor{critred}{\textbf{FAIL}} \\
Mathematics/Equations & Agent 4 & 51/100 & $\sim$80/100 & \textcolor{mingreen}{\textbf{PASS (post-fix)}} \\
Tables/Documentation & Agent 5 & 75/100 & $\sim$85/100 & \textcolor{majblue}{\textbf{Adj.$R^2$ pending}} \\
JF Referee & Agent 6 & 60/100 & --- & \textcolor{critred}{\textbf{FAIL for JF}} \\
& & & & \textcolor{majblue}{\textbf{$\sim$75 for JFQA}} \\
\midrule
\textbf{Weighted Aggregate} & & \textbf{$\sim$71} & \textbf{$\sim$80} & \textcolor{majblue}{\textbf{JFQA threshold}} \\
\bottomrule
\end{longtable}
\end{center}

\noindent\textbf{Weight schema} (renormalized; code/lit/data not yet scored):\\[0.3em]
Manuscript polish: 15\% $\cdot$ Consistency: 15\% $\cdot$ Claims: 20\% $\cdot$
Equations: 15\% $\cdot$ Tables: 15\% $\cdot$ Referee: 20\%

\vspace{1em}
\noindent\textcolor{critred}{\textbf{Agent 3 (Unsupported Claims) blocks submission}} until the welfare language
in the conclusion is resolved (add return evidence or remove claim) and "consistent only with"
is weakened. These are text edits, not data work.

\vspace{1em}
\noindent\textbf{Path to JF:}
\begin{enumerate}[leftmargin=2em]
  \item Complete Section~3 (IV, Wald test, returns, autocorrelation diagnostic, remove placeholders).
  \item Complete Section~4 (Adj.~$R^2$, welfare language, missing citations, table fixes).
  \item Re-run all six review agents. Target aggregate $\geq 85$.
  \item Submit to JF with full replication package.
\end{enumerate}

\vspace{1em}
\noindent\textbf{Path to JFQA/RFS (shorter):}
\begin{enumerate}[leftmargin=2em]
  \item Complete Section~2 completed fixes (done).
  \item Complete Section~4 priority items (2--4 weeks).
  \item Remove "we plan to" language (Section~3.5, text edit, 1 hour).
  \item Submit.
\end{enumerate}

\end{document}
